\chapter{Peripheral Nerve Disorders}
\index{Peripheral Nerve Disorders}

\medskip
\noindent\textit{This chapter is for general education. New weakness, rapidly spreading numbness, or breathing problems need urgent assessment.}

\section*{What it is}
Peripheral nerve disorders affect nerves outside the brain and spinal cord. They may involve sensory nerves, motor nerves, or the autonomic nerves that control functions such as sweating, blood pressure, digestion, and bladder function. Causes include diabetes, compression or injury, nutritional deficiency, autoimmune disease, inherited conditions, infections, medicines, toxins, and sometimes no identified cause.

\section*{Symptoms}
Symptoms can include numbness, tingling, burning or electric-shock pain, sensitivity to touch, weakness, cramps, loss of balance, reduced reflexes, dizziness on standing, abnormal sweating, or bowel and bladder changes. The pattern and speed of progression help guide assessment.

\section*{Diagnosis}
Assessment includes a history, examination of sensation, strength, reflexes, and gait, and review of medicines and exposures. Blood tests, nerve-conduction studies, electromyography, imaging, or specialist review may be needed. Finding and treating the cause is often more important than treating the nerve symptoms alone.

\section*{Treatment and daily care}
Treatment may include managing diabetes or another cause, prescribed medicines for nerve pain, physiotherapy, strength and balance work, occupational therapy, and foot protection. Check feet and skin regularly when sensation is reduced, use safe footwear, and avoid heat sources that cannot be felt. Do not stop a suspected medicine without medical advice.

\section*{When to Seek Medical Care}
Seek emergency help for sudden weakness or numbness, facial droop, trouble speaking, loss of bladder or bowel control, rapidly ascending weakness, difficulty breathing or swallowing, severe dizziness, or new inability to walk. Arrange review for progressive symptoms, unexplained falls, wounds, or persistent pain.

\section*{Sources and further reading}
\begin{itemize}
\item National Institute of Neurological Disorders and Stroke. \textit{Peripheral neuropathy}. \url{https://www.ninds.nih.gov/health-information/disorders/peripheral-neuropathy}
\item MedlinePlus. \textit{Peripheral neuropathy}. \url{https://medlineplus.gov/peripheralneuropathy.html}
\end{itemize}
